\subsection{Two 1990 Origins: Dyna and Differentiable World Models}
\label{section:fc_origins1990}

Two papers published in 1990 introduced what the modern literature now calls a world model. In each, an agent learns a forecaster of its environment from interaction data and uses simulated experience drawn from that forecaster to improve its policy, rather than waiting for real transitions to accumulate. \citet{sutton1990} arrived at the architecture through dynamic programming and reinforcement learning, framing it as the unification of direct and indirect updates within a single control loop. \citet{Schmidhuber1990} arrived at it through recurrent neural networks, framing it as a way to make the environment itself differentiable for credit assignment. The two papers do not cite one another and use different vocabularies, but the architectural commitment is the same. We treat them in parallel rather than in lineage, since each anticipates a distinct branch of the modern literature.

\subsubsection{Sutton's integrated architecture}
\label{section:fc_origins_sutton}

\citet{sutton1990} proposes Dyna as an architecture that interleaves three processes at every real time step. The agent performs a direct Q-learning update from the observed transition $(s, a, r, s')$, then updates a learned model $\widehat{P}(s' \mid s, a)$ and $\widehat{r}(s, a)$ from the same transition, then performs $K$ planning updates by sampling previously visited state-action pairs $(\tilde{s}, \tilde{a})$, querying the model for simulated outcomes $(\tilde{r}, \tilde{s}')$, and applying the same Q-update to those synthetic transitions. The ratio $K$ of planning steps to real steps is the lever that controls how heavily the agent exploits its learned model between actions, and is the first explicit statement of the sample-efficiency claim that modern model-based work continues to revisit; the MBPO short-rollout analysis in \citet{janner2019model} is in the same family. We defer the algorithm block to \S\ref{section:fc_dyna_q}, where we develop Dyna-Q in the notation of this survey.

The architectural openness of Dyna is as important as its sample-efficiency claim. Sutton states explicitly that the learned model can be probabilistic and oftentimes incorrect, with the planner doing useful work as long as the cost of synthetic experience is small relative to the cost of real experience. No structure is imposed on the class of admissible models, on the loss used to fit them, or on the planner that consumes their output. This is the architectural slot that subsequent work fills with tabular counts, Gaussian processes, deep ensembles, recurrent state-space models, and implicit latent dynamics, each plugged into the same outer loop. The conceptual shift from a fixed control algorithm to an open architecture with a model component is what makes Dyna the entry point for the line that runs through \S\ref{section:fc_mbpo} and \S\ref{section:fc_tdmpc2}.

\subsubsection{Schmidhuber's controller-model architecture}
\label{section:fc_origins_schmidhuber}

\citet{Schmidhuber1990} proposes a paired architecture consisting of a recurrent model network $M$ and a controller network $C$. The model $M$ receives the controller's outputs together with current observations and is trained, self-supervised, to predict the next inputs to $C$, including sensory observations and reinforcement signals. The controller $C$ produces actions and is trained jointly with $M$. Once $M$ is a reasonable forecaster, the environment becomes differentiable from the controller's perspective, since the Jacobian of $M$ stands in for the unknown Jacobian of the true dynamics. Planning is then implemented as gradient descent through the model network, by unrolling $C$ and $M$ together for a fixed horizon and descending on a predicted pain or cost signal, without executing any real actions during the planning phase. This is the first proposal of analytic gradient backpropagation through a learned dynamics model for policy improvement.

In the notation of this survey, let $M_\theta$ parameterize the world model with $p_\theta(s_{t+1} \mid s_t, a_t)$ and let $C_\phi$ parameterize the controller with $\pi_\phi(a_t \mid s_t)$. The model is fit on observed transitions by the negative log-likelihood
\[
  \mathcal{L}_M(\theta) = -\sum_t \log p_\theta(s_{t+1} \mid s_t, a_t),
\]
and the controller is fit by maximizing the expected discounted return under the \emph{learned} model,
\[
  J(\phi) = \mathbb{E}_{a_t \sim \pi_\phi(\cdot \mid s_t),\; s_{t+1} \sim p_\theta(\cdot \mid s_t, a_t)} \Big[ \sum_t \gamma^t r_t \Big],
\]
with the planning operator $\phi \leftarrow \phi + \eta \, \nabla_\phi J(\phi)$ implemented by backpropagating through the $H$-step unrolled $M_\theta$ via reparameterized or score-function gradients. The curiosity unit adds an intrinsic reward $r_t^{\text{int}} = \| s_{t+1} - \mathbb{E}_{M_\theta}[s_{t+1} \mid s_t, a_t] \|^2$ to the extrinsic reward inside $J$, with $M_\theta$ trained to drive precisely this prediction error toward zero so that regions of state space where $M_\theta$ is currently uncertain attract $C_\phi$. The idea reappears in the deep era as the intrinsic curiosity of \citet{Pathak2017}, instantiated with a forward-dynamics network in a learned feature space and used to drive exploration in sparse-reward video games; the full deep realization of the controller-model architecture is the work of \citet{HaSchmidhuber2018} and is treated in \S\ref{section:fc_ha_schmidhuber}.

The contrast with Sutton's Dyna is sharp on two dimensions and shallow on a third. The model class differs, a deterministic tabular Dirac in Dyna versus a differentiable neural network in Schmidhuber. The planner differs, $K$ tabular $Q$-updates per real step in Dyna versus a gradient step on $\phi$ through the unrolled $M_\theta$ in Schmidhuber. The outer loop is the same: learn $M$ from real data, use $M$ to improve the policy, act in the world.

The symmetry between the two 1990 papers is what makes them an instructive joint origin. Both describe an agent that learns a forecaster of its environment and uses simulated experience from that forecaster to improve its policy, and both do so without the empirical tools that would let the idea work at the scale of pixel inputs or continuous control benchmarks. Both anticipate components that took roughly three decades to mature into reliable practice, intrinsic curiosity in one direction and planning amplification through short model rollouts in the other. In the next subsection we develop Dyna-Q in the notation of this survey and state the empirical claim that the planning-step parameter $K$ trades real samples for synthetic ones, which is the cleanest formal target the 1990 papers leave for the chapters that follow.
